% Sources: experiments/v326_two_action_upper/{RESULTS.md,PROTOCOL.md,config.json,review01/aggregate.json}
% Sources: experiments/v325_budget_permission/{RESULTS.md,review01/aggregate.json}
% Sources: experiments/v316_upper_no_cv/learning01/{RESULTS.md,config.json}
% Sources: experiments/v331_candidate_capture/RESULTS.md
% Source: experiments/v350_live_basis_wiring/matched_launch01/review01/matched_aggregate.json
% Source: experiments/v352_budget_interface/review01/aggregate.json
% Source: experiments/v399_equal_budget_manual/review01/compare.json
% All estimates below are transcribed from existing server aggregates, not recomputed.

\section{Experimental Evaluation}
\label{sec:experiments}

The evaluation separates three questions: how learned and strong manual joint
control compare at matched budget retention, how retention changes service and
regularity, and what harmonic terms contribute under matched retraining.
Historical objective and authorization comparisons provide additional context.
The preselected default candidate and its retention of 0.5 remain unchanged.

\subsection{Experimental Design and Comparators}

The simulated corridor has 22 physical stops, including two terminals, with
21 segments of 500 m in each direction (10.5 km one way). Each bus has capacity
50. The public timetable specifies 131 trips per direction at 360-s intervals;
the two initial departures are offset by 180 s. Time zero corresponds to
06:00. The simulation advances in 1-s steps, updates road conditions every
300 s, and generates passengers in 20-s windows from the configured hourly
origin--destination table. Counts are Poisson, with arrival times sampled
uniformly within each window and revealed only when they occur. The scenario
generator perturbs hourly demand and shifts the hourly template. Link speeds
are sampled around hourly means, clipped to the supported range, and used in
the physical vehicle update. This is a stochastic simulation study; the
development demand seeds are not independent observations from an operational
deployment.

The nominal demand cutoff is the last public departure at 46,980 s. A
14,400-s service tail gives a maximum accounting horizon of 61,380 s; physical
driving can end earlier when service has finished. Passenger metrics are
restricted accumulated times over this evaluation horizon. They include
waiting for generated passengers who remain unserved, rather than conditioning
on completed journeys. Accordingly, ``journey'' below means the sum of
accounted waiting and onboard exposure per generated passenger, not an assertion
that every generated passenger has reached a destination.

All reported full-day comparisons retain the same 12-vehicle fleet, 262 scheduled
trips, demand realization within each paired day, and fast holding controller.
Passenger accounting includes the prescribed waiting and onboard tail. Waiting
and onboard time are divided by the number of passengers generated on the same
day; journey time is their sum. Vehicle-readiness delay is normalized by 262
trips, while public departure delay is reported as a daily total relative to the
unchanged public timetable. Consequently, revising an internal operating plan
cannot erase a public delay by redefining its reference. Actual headway
\(CV^2\) is reported separately from these service measures.

The four-mode learner and FreqDuet each use three initialization seeds, 314011,
314037, and 314049. Each seed receives 32 distinct training demand days, identifiers
1780001--1780032, in eight batches of four days. A batch uses its preceding
checkpoint and is followed by 192 updates, giving 1,536 updates per seed and
4,608 in total. The final checkpoint is evaluated deterministically without
updates on demand days 2077001--2077004. Seed 314011 and its eighth-batch
checkpoint were designated before evaluation as the primary candidate; the
other seeds assess sensitivity rather than provide alternative candidates for
selection. The four evaluation days had already been used during development.

The comparators isolate different parts of the coordination mechanism.
\emph{Fixed plan} retains the original upper timetable with the same fast
controller. The \emph{four-mode learner} chooses among fixed, plan, budget, and
joint permissions using the passenger-and-delay objective. The
\emph{permission-mapped learner} freezes that actor and maps fixed to budget and
plan to joint, retaining its weights, prediction inputs, and fast controller.
This intervention provides budget authorization at every upper decision but
does not retrain the critic under the changed action semantics.
\emph{Manual joint} always authorizes the same low-frequency plan rule and
holding-budget mechanism. It is therefore a strong coordinated comparator,
not an uncontrolled schedule. FreqDuet instead learns directly between the two
executable budget and joint actions, with actor, critic, replay labels, previous
action encoding, and execution sharing the same semantics (archived as V326,
V316, and V325 for FreqDuet, the four-mode learner, and the permission-mapped
learner, respectively). For seed-sensitivity comparisons, the four-mode and
permission-mapped references both use their frozen seed-314011 policies; manual
joint is also shared across the three FreqDuet seeds.

The optimization objective is given in~\eqref{eq:objective}.

For the main comparison, both learned and manual joint control use retention
$\rho_B=0.5$ and $0.25$. The latter manual arm is evaluated on the same four
days with no learning; the learned arms and the default manual arm are reused
from the existing evaluations. Equal retention means the same lower bound in
the allowance formula, not equal realized allowance, holding, or state
trajectory. The additional manual paths each complete all 262 trips, preserve
the 12-vehicle pool, and match 5,200 common random events per day. The original
default manual path is reproduced before applying the tighter setting.

Differences are always method minus the named comparator, so negative values
denote reductions. The archived summaries use equal day weights and all $4^4=256$ paired
resamples with replacement from the four days. We retain their descriptive 95\% intervals
and report each learning seed separately. Twelve evaluation runs represent
three policies on four demand days, not twelve independent demand days.

\subsection{Main Comparison at Matched Budget Retention}
\label{sec:equal_budget}

Table~\ref{tab:equal_budget} reports the four controller--retention combinations.
At retention 0.25, the learned policy reduces mean journey time by only
0.094~s/passenger relative to strong manual joint control. The corresponding
waiting, journey, readiness, and actual $CV^2$ intervals all span zero.
Journey differences by day are $+0.883$, $+0.163$, $+0.336$, and $-1.758$
s/passenger; waiting also improves only on the fourth day. At default retention
0.5, the mean journey difference is larger at $-1.925$~s/passenger, but its
interval also spans zero. The main outcome is therefore a matched description
of operational performance, not an established learning advantage.

\begin{table*}[t]
\centering
\caption{Matched-retention main comparison on four reused development days.
The learned policy is the preselected K4 candidate; manual joint is a strong
coordinated controller. Retention is an allowance floor, not realized holding.
Intervals in panel (b) are descriptive paired 95\% intervals.}
\label{tab:equal_budget}
\small
\textbf{(a) Absolute means}\par\smallskip
\setlength{\tabcolsep}{8pt}
\begin{tabular}{lrrrrr}
\toprule
Upper policy & $\rho_B$ & Waiting (s/passenger) & Journey (s/passenger) & Readiness (s/trip) & Actual $CV^2$ \\
\midrule
Manual joint & 0.50 & 308.775 & 998.705 & 65.450 & 0.08083 \\
Learned K4 (default) & 0.50 & 307.752 & 996.780 & 64.725 & 0.07802 \\
Manual joint & 0.25 & 307.372 & 995.285 & 62.732 & 0.08273 \\
Learned K4 & 0.25 & 307.300 & 995.191 & 62.179 & 0.08220 \\
\bottomrule
\end{tabular}
\par\medskip
\textbf{(b) Learned K4 minus manual joint at the same retention}\par\smallskip
\begin{tabular}{lrr}
\toprule
Metric & $\rho_B=0.50$ & $\rho_B=0.25$ \\
\midrule
Waiting (s/passenger) & $-1.023\;[-3.396,0.668]$ & $-0.072\;[-1.307,0.955]$ \\
Journey (s/passenger) & $-1.925\;[-5.935,0.762]$ & $-0.094\;[-1.234,0.703]$ \\
Readiness (s/trip) & $-0.725\;[-3.295,0.885]$ & $-0.552\;[-1.552,0.447]$ \\
Actual $CV^2$ & $-0.002816\;[-0.007316,0.000185]$ & $-0.000532\;[-0.001479,0.000620]$ \\
\bottomrule
\end{tabular}
\end{table*}


These small aggregate differences do not arise from identical execution. At
retention 0.25, the controllers differ in 43 trips' published release times,
245 trips' actual departures, and 3,989 matched holding commands across the
four days. These are records from different closed-loop trajectories; they
cannot identify the passenger effect of any individual command.

The matched comparison also clarifies the earlier cross-budget contrast.
Learned control at 0.25 versus manual joint at 0.5 reduces journey time by
3.514~s/passenger. Along the comparison sequence that first tightens the
manual budget and then changes upper policy at retention 0.25, this difference
is $-3.420-0.094=-3.514$~s/passenger. The waiting contrast similarly decomposes
as $-1.403-0.072=-1.475$~s/passenger. These are accounting identities along a
stated sequence, not a unique allocation of interacting causal effects.
The journey difference-in-differences is $+1.831$~s/passenger
$[-2.039,6.681]$; there is no resolved advantage in how the learned policy
responds to budget tightening. Section~\ref{sec:budget_settings} reports the
within-policy configuration effects.

\subsection{Historical Objective and Authorization Comparisons}

The historical comparisons evaluate the default preselected candidate
against earlier learners and manual joint, all at retention 0.5.
Relative to the four-mode learner, journey
time decreases by 7.448~s/passenger, vehicle-readiness delay by
12.221~s/trip, and public departure delay by 3,162.5~s/day. All four days improve
on these three measures, and their descriptive intervals remain below zero.
The largest passenger component is onboard time, reduced by
5.873~s/passenger. Holding decreases by 14.704~s/trip. The waiting reduction is
smaller and varies across days, indicating that the clearest recovered benefit
is reduced onboard burden and improved vehicle circulation.

The comparison with manual joint is materially closer. FreqDuet reduces mean
waiting by 1.023~s/passenger and journey time by 1.925~s/passenger, with lower
mean readiness and public delay. However, the corresponding intervals cross
zero. Waiting increases by 0.315 and 1.022~s/passenger on the third and fourth
days; journey time increases by 1.308~s/passenger on the fourth day. The
aggregate result therefore identifies a competitive candidate with favorable
means on this panel, rather than a demonstrated uniform improvement over
manual coordination. The comparator is retained in the main table because the
larger improvement over the four-mode learner alone would overstate the
remaining benefit attributable to learning.

Against the permission-mapped learner, retraining adds only
0.365~s/passenger of mean journey reduction and 0.271~s/passenger of waiting
reduction. Readiness and public delay increase slightly. This small difference
shows that most of the full-day recovery was already present after enabling
budget authorization. Correcting the action representation makes learning
consistent with execution, but does not itself establish an additional traffic
benefit.

\subsection{Objective and Authorization Ablations}

The objective ablation retains the four-mode architecture and training budget
while removing the upper regularity term from
\(J_1=J_0+CV^2/0.05\). For the preselected seed, the resulting four-mode learner
reduces journey time by 7.258~s/passenger relative to its regularity-penalized
counterpart. Relative to fixed plan, its journey reduction is
7.047~s/passenger [\(-11.169,-2.360\)], readiness decreases by
23.120~s/trip, and public delay decreases by 6,007~s/day. The two other seeds
also reduce journey time against fixed plan, by 12.296 and 9.987~s/passenger.
The action records show that budget and joint permissions become active after
the objective change; the regularity-penalized learners had selected neither
mode in their final evaluations.

This recovery includes a real regularity tradeoff. For the preselected seed,
actual \(CV^2\) increases by 0.01894 [0.01334,0.02454] against fixed plan.
Although \(J_0\) decreases by 0.10374, evaluating the same trajectories with
\(J_1\) increases cost by 0.27498. Thus the objective change exposes a useful
passenger--circulation tradeoff; it does not make the old composite criterion
pass. The four-mode learner also remains 5.523~s/passenger worse than manual
joint in journey time, on average, with all four days worse.

The authorization ablation then freezes this actor and changes only its
permission mapping. Relative to the unmapped four-mode learner, it reduces
journey time by 7.083~s/passenger [\(-15.617,-1.896\)], readiness by
12.658~s/trip [\(-22.407,-4.824\)], and public delay by
3,324.5~s/day [\(-5883.8,-1275.5\)]. Its onboard reduction is
5.779~s/passenger, while the waiting reduction of 1.304~s/passenger has an
interval spanning zero. Persistent-budget accounting confirms that existing
allowances are consumed rather than replenished at repeated decision times.
These results support authorization as an important source of the recovered
benefit. They do not isolate every indirect effect, because changed holding
also changes subsequent vehicle supply and plan opportunities.

\subsection{Learning Sensitivity and Executed Plan Authority}

Table~\ref{tab:seed_results} reports all final learning seeds. Each has lower
mean waiting and journey time than manual joint, but only the preselected seed
improves mean waiting against the permission-mapped learner. The other seeds
increase waiting by 0.341 and 0.270~s/passenger against that comparator. The
budget/joint action counts also differ substantially, despite identical action
semantics and training budgets. This variation limits the interpretation of the
small incremental gain from retraining.

\begin{table}[t]
\centering
\caption{Seed sensitivity. Passenger differences are in s/passenger. MJ denotes
manual joint; PM denotes the permission-mapped learner.}
\label{tab:seed_results}
\setlength{\tabcolsep}{3pt}
\footnotesize
\begin{tabular}{lrrrr}
\toprule
Seed & $\Delta W$/MJ & $\Delta(W+R)$/MJ & $\Delta W$/PM & Budget/joint \\
\midrule
314011 & $-1.023$ & $-1.925$ & $-0.271$ & 70/118 \\
314037 & $-0.411$ & $-0.942$ & $0.341$ & 174/14 \\
314049 & $-0.482$ & $-0.822$ & $0.270$ & 106/82 \\
\bottomrule
\end{tabular}
\end{table}

The upper controller is operationally active. Under the archived criterion
requiring a matched publication, bound dispatch, and at least a 5-s actual
departure difference, the three seeds have material plan effects on four, two,
and four demand days. The preselected candidate affects seven, five, six, and
one trips across the four days. Its actor disagrees with its online twin-min
critic ranking on 51 of 188 decisions; the individual critics disagree on
38. These diagnostics distinguish valid execution from unresolved value and
policy disagreement. A critic preference is not a measured counterfactual
traffic outcome, so these counts cannot be converted into lost passenger
seconds or used to justify replacing the evaluated actor.

\subsection{Matched Harmonic-Basis Ablation}
\label{sec:harmonic_ablation}

We retain the evaluated $K=4$ harmonic representation and retrain a $K=0$
comparison from initialization with identical seeds, training days, update
counts, reward, fast controller, and action semantics. Both the actor input and
timetable law consume the same causal summary object. A frozen complete day verifies $K=4$ operational reproduction before this
ablation. The $K=0$ arm retains the intercept, trend, online adaptation, and
innovation statistics; removing four sine--cosine pairs reduces estimator
capacity from ten to two coefficients. Thus this is a matched harmonic-basis
ablation, not removal of every temporal summary.

\begin{table}[t]
\centering
\caption{Matched retraining ablation: $K4-K0$ for the preselected candidate.
Negative differences favor retaining harmonics.}
\label{tab:harmonic}
\setlength{\tabcolsep}{4pt}
\footnotesize
\begin{tabular}{lr}
\toprule
Metric & Mean [descriptive 95\% interval] \\
\midrule
Waiting (s/passenger) & $-0.404\;[-1.352,0.316]$ \\
Onboard (s/passenger) & $-0.657\;[-0.951,-0.387]$ \\
Journey (s/passenger) & $-1.061\;[-1.889,-0.233]$ \\
Readiness (s/trip) & $-0.441\;[-0.955,0.049]$ \\
Public delay (s/day) & $-93.75\;[-228.50,21.00]$ \\
$J_0$ & $-0.00713\;[-0.01249,-0.00178]$ \\
Actual $CV^2$ & $-0.000824\;[-0.002377,0.000730]$ \\
\bottomrule
\end{tabular}
\end{table}

The preselected $K4$ candidate reduces journey time on all four days
(Table~\ref{tab:harmonic}). Mean journey differences for seeds 314011, 314037,
and 314049 are $-1.061$, $-0.286$, and $-1.174$ s/passenger, respectively;
9 of 12 seed--day comparisons favor $K4$. Onboard time improves in 10 of 12
comparisons, whereas waiting improves in only six. The result supports a
modest journey benefit from retaining harmonic structure under this training
recipe. It does not establish a stable waiting benefit. Because the
intervention changes both representation capacity and the summary supplied to
the actor and plan law, their individual contributions remain coupled.
The main candidate and historical comparison rules are unchanged.

\subsection{Prespecified Budget Configuration}
\label{sec:budget_settings}

We evaluate $\rho_B\in\{0.25,0.5,0.75\}$ in~\eqref{eq:budget} using the
frozen preselected $K4$ policy. The default remains 0.5. Before evaluating the
two other settings on the same four days, one default full day reproduces all
compared actions, event times, policy states, rewards, probabilities, and
service outcomes exactly. All nine paths finish without retraining or retries;
physical, budget, action, reward, and paired random-event accounting pass.
The settings were declared in advance, and no setting replaces the main policy.

\begin{table}[t]
\centering
\caption{Prespecified budget settings relative to default $\rho_B=0.5$.
Four-day means; negative costs improve. The default contrast is zero.}
\label{tab:budget_settings}
\setlength{\tabcolsep}{6pt}
\footnotesize
\begin{tabular}{lrr}
\toprule
Metric & $\rho_B=0.25$ & $\rho_B=0.75$ \\
\midrule
Waiting (s/passenger) & $-0.452$ & $+2.015$ \\
Onboard (s/passenger) & $-1.137$ & $+2.916$ \\
Journey (s/passenger) & $-1.589$ & $+4.931$ \\
Readiness (s/trip) & $-2.546$ & $+4.401$ \\
Public delay (s/day) & $-644.75$ & $+1153.25$ \\
Holding (s/trip) & $-0.969$ & $+5.940$ \\
Actual $CV^2$ & $+0.004177$ & $+0.000196$ \\
\bottomrule
\end{tabular}
\end{table}

At 0.25, journey and readiness improve on all four days, with descriptive
95\% intervals $[-2.641,-0.541]$ s/passenger and $[-2.985,-2.168]$ s/trip.
Waiting improves on only two days; its interval is $[-1.197,0.293]$
s/passenger. Actual $CV^2$ increases on every day, with interval
$[0.000919,0.008926]$. At 0.75, journey and readiness worsen on every day,
and the waiting, journey, readiness, and public-delay intervals are above zero.
Table~\ref{tab:budget_settings} thus identifies a service--regularity tradeoff,
not a uniformly superior setting.

The tighter setting changes 4,662 matched holding actions, 30 published
releases, and 345 actual departures across the panel. Joint permissions
increase from 118 to 122 despite frozen weights: changed budgets alter later
states and hence policy decisions. Executed holding falls by 253.75 s/day;
reductions within budgeted periods are partly offset elsewhere on the changed
trajectories. These records establish a physical effect beyond changing a
logged allowance.

Against manual joint at its original retention 0.5, retention 0.25 reduces
waiting by 1.475 s/passenger $[-3.419,-0.206]$, journey by 3.514 s/passenger
$[-7.095,-0.760]$, and readiness by 3.271 s/trip $[-5.898,-1.635]$.
Waiting improves on three days and increases by 0.013 s/passenger on the fourth;
journey and readiness improve on all four. This comparison changes both
budget configuration and upper control, so it does not isolate a learning
increment at matched retention. Conversely, retention 0.75 improves $CV^2$
against manual joint on every day while worsening journey and readiness.

\subsection{Direct Coupling Between Timetable and Holding}

A paired replay provides a concrete mechanism example beyond aggregate plan
activity. At one recorded query, the effective vehicle, queue, forecast,
history, and holding-budget inputs are identical across two branches. Only the
published plan differs: one trip's release and target move by $+1$~s, and
another target moves by $-1$~s. The native reference action is 30~s, the remaining
budget is 66~s, and the same seven action bins are available.

\begin{table}[t]
\centering
\caption{Same-state timetable-to-holding interaction. Costs are predicted
passenger-seconds; a positive difference favors 20~s over 15~s.}
\label{tab:coupling}
\setlength{\tabcolsep}{4pt}
\footnotesize
\begin{tabularx}{\columnwidth}{lrrX}
\toprule
Plan branch & $C(15)-C(20)$ & Selected & Rejected alternative \\
\midrule
Budget & $142.345$ & 20~s & 15~s: 2-s plan-relative delay violation \\
Joint & $204.892$ & 15~s & 20~s: up to 1-s plan-relative delay violation \\
\bottomrule
\end{tabularx}
\end{table}

Both branches rank 20~s ahead of 15~s by passenger cost, but their plan-relative
feasible sets differ (Table~\ref{tab:coupling}). The selected action changes
from 20 to 15~s because the published plan changes which downstream departure
constraints bind. This establishes an implemented information path from the
slow plan to fast holding, rather than two disconnected controllers. It is a
local mechanism result, not a full-day synergy estimate. At another captured
query, vehicle and demand inputs have already diverged, so its changed action
cannot be attributed to the plan alone.

The same replay also illustrates the objective's meaning. The selected joint
action, relative to its 30-s reference, predicts 3,623.636 additional
person-seconds of waiting and 3,660.720 fewer onboard person-seconds, a net
37.083-person-second reduction. Minimizing their sum can therefore select an
action that increases waiting. Separating these components is essential when
interpreting a journey-time improvement.
